\cotefile{156-8}
\scanpage{34}
\mspage{25}
% archive sheet number: 34
% inv: text | top | Magasins triviaux
\uline{Magasins triviaux}\quad Def : Un magasin est dit ``trivial'' si $X\overset{\circ}{\ll}Y\Longleftrightarrow X=Y$, et si $X\mathrel{|\circ|}Y\Longleftrightarrow X\ne Y$, donc la seule donnée qui reste à déterminer est $X\leq Y$. Mais on \unsure{vérifie} \inserted{trivialement} que pour \deleted{tout} ens. ordonné $(M,\leq)$, si on prend \deleted{$X\overset{\circ}{\ll}Y$}
% inv: display | upper third | X circle-ll Y
\[
X\overset{\circ}{\ll}Y\overset{\text{déf}}{\Longleftrightarrow}X=Y,\qquad X\mathrel{|\circ|}Y\overset{\text{déf}}{\Longleftrightarrow}X\ne Y,
\]
% inv: text | upper middle | on trouve une structure
on trouve une structure de magasin.

% inv: deletion | upper middle | Ce magasin ne satisfait
\deleted{Ce magasin ne satisfait Mag L 1 que si \unsure{tout}} \deleted{\unsure{fig.}} \deleted{$L=$ ens des \unsure{éléts} minimaux}

% inv: display | middle | X ll Y iff X le Y
\[
\begin{aligned}
X\ll Y&\Longleftrightarrow X\leq Y\\
X\unsure{\gtrless|\bowtie}Y&\qquad\text{tjrs satisfait}\\
X\parallel Y&\qquad\text{\unsure{,,,}}\quad\widetilde X\cap\widetilde Y=\varnothing\quad\text{i.e.}\quad\{X,Y\}\text{ pas minoré}.
\end{aligned}
\]
% inv: display | middle | L equals ens des éléments minimaux
\[
L=\text{ens des éléments minimaux de }X
\]
% inv: display | middle lower | omb(X)
\[
\begin{aligned}
\mathrm{omb}(X)&=\{a\in L\mid a\leq X\}\\
\mathrm{omb}(X)^\circ&=\deleted{\{a\in L\mid a<X\}}\begin{cases}\varnothing\text{ si }X\notin L\\\mathrm{omb}(X)=\{X\}\text{ si }X\in L\end{cases}
\end{aligned}
\]

% inv: margin note | slanted left upper | NB Conditions équivalentes
% margin note written sideways
% slanted
\margin{\uline{NB} Conditions \unsure{équivalentes} pour un magasin $M$ :\par
(i) $M$ trivial\par
(ii) $X\gtrless Y$ tjrs satisfait\par
(iii) $\unsure{\mathrm{FigMag}|\mathrm{Fig}_{\mathrm{Mag}}}(M)=$ ens. des figures \unsure{initiales|filtrantes} pour $\leq$\par
(iv) $M\in\mathrm{Fig}_{\mathrm{mag}}(M)$\par
(v) $X\ne Y\Longrightarrow X\mathrel{|\circ|}Y$}
% inv: margin note | slanted left middle | Ces conditions impliquent
% margin note written sideways
% slanted
\margin{Ces conditions impliquent les conditions suivantes qui \unsure{sont} \unsure{équivalentes}, mais qui n'impliquent pas les précédentes\par
(vi) $X\overset{\circ}{\ll}Y\Longrightarrow X=Y$\par
(vii) $X\ll Y\Longrightarrow X\leq\deleted{\unsure{\gtrless}}Y$}

% inv: text | lower middle | La condition Mag L 1
La condition \uuline{Mag L 1} i.e. $\mathrm{omb}(X)^\circ\ne\unsure{0|\varnothing}\ \forall X$, \unsure{signifie} que tout $X$ est minimal i.e. la relation d'ordre est discrète. On voit qu'un magasin trivial ne satisfait \unsure{quasiment|pratiquement} jamais Mag L 1, et encore moins Mag L 2 --- sauf s'il est doublement trivial $(X\leq Y\Longleftrightarrow X=Y)$.

% inv: margin note | slanted left lower | Les magasins triviaux
% margin note written sideways
% slanted
\margin{Les magasins triviaux jouent le rôle d'ensembles d'indices pour des familles de magasins \unsure{généraux}, ou des \uline{maquettes} \unsure{combinatoires} de figures. On peut \unsure{appeler} \unsure{simplement} un tel magasin une \uline{maquette} (!), (une \unsure{m.} \unsure{squelette} (\unsure{catégorique})).}

% inv: text | bottom | Figure indexée
\uline{Figure indexée} par un ens. ordonné $I$ : figure $F$, \unsure{munie} d'un isom.
% inv: diagram | bottom | I isomorphic to tilde F
\diagram{\begin{tikzcd}[ampersand replacement=\&]
\deleted{\unsure{\in}}\quad I \arrow[r,"\sim",dash] \& \widetilde F\\
i \arrow[r,mapsto] \& X_i
\end{tikzcd}}
